\documentclass[11pt,a4paper]{article}
\usepackage[margin=25mm]{geometry}
\usepackage[T1]{fontenc}
\usepackage{lmodern,microtype,amsmath,amssymb,booktabs,graphicx,xcolor,array}
\usepackage[colorlinks=true,allcolors=blue!45!black]{hyperref}
\usepackage{fancyhdr}
\makeatletter\let\inputtablerows\@@input\makeatother
\pagestyle{fancy}\fancyhf{}
\lhead{\small Constitutive complexity and calibration information}
\rhead{\small Development manuscript}
\cfoot{\thepage}\setlength{\headheight}{14pt}
\setlength{\parskip}{5pt}\setlength{\parindent}{0pt}
\IfFileExists{generated/paper_inputs.tex}{\input{generated/paper_inputs.tex}}{
\newcommand{\AnnotatedFormulationCount}{pending}
\newcommand{\PilotModelCount}{pending}
\newcommand{\PilotEvaluationTestCount}{pending}
\newcommand{\PilotMaximumStrainPercent}{pending}
\newcommand{\PilotCalibrationSamples}{pending}
\newcommand{\PilotVolumeScale}{pending}
\newcommand{\PilotStartCount}{pending}
\newcommand{\CompressionReplayCount}{pending}
\newcommand{\CompressionStartThresholdCount}{pending}
\newcommand{\CompressionStartThresholdKPa}{pending}
\newcommand{\CompressionObservedStartKPa}{pending}
\newcommand{\CompressionCoarseIncrement}{pending}
\newcommand{\CompressionFineIncrement}{pending}
\newcommand{\CompressionStressRatios}{pending}}
\providecommand{\SurveyOccurrenceCount}{pending}
\providecommand{\SurveyPublicationCount}{pending}
\providecommand{\SurveyTitleDecisionCount}{pending}
\providecommand{\SurveyUntitledCount}{pending}
\providecommand{\SurveyRetrieveCount}{pending}
\providecommand{\SurveyBackgroundCount}{pending}
\providecommand{\SurveyUnresolvedCount}{pending}
\providecommand{\SurveyOutsideScopeCount}{pending}
\title{\textbf{What does constitutive complexity buy?}\\[5pt]
\large Predictive capability and calibration requirements in geomechanics\\
\normalsize Study methods and development evidence}
\author{Benjy Marks \and Fran\c{c}ois Guillard}
\date{31 October 2026 --- working manuscript, not a submission}
\begin{document}
\maketitle
\begin{abstract}
Constitutive models differ in the mechanisms they represent, the information needed to initialize and calibrate them, and the predictions they can support. A useful comparison must connect these differences to specific experimental tasks. This manuscript develops a source-traceable model survey and a controlled material-point comparison addressing predictive benefit, calibration-information requirements and the value of model extensions. The planned broad survey targets a future collection of more than fifty formulations; the current annotations and development experiments do not establish that coverage. Completed pilot work exercises whole-test calibration splits, pressure and density transfer, physical error scales, failure accounting, numerical refinement and optimization sensitivity on Karlsruhe fine sand. A matched-count development comparison examines three pressure-spanning versus three density-spanning calibration tests within each formulation. A separate fixed-parameter compression diagnostic finds initialization and strain-coordinate sensitivities larger than the observed numerical-refinement differences. These development results expose information and convention requirements; supplied prior constants, model review status and experimental metadata still prevent a final model ranking or predictive validation.
\end{abstract}

\textbf{Status and scope.} The study freeze, broad survey, independently reviewed benchmark admission and final evaluation remain open. Values below are generated from source-pinned development records. The original software-focused draft is retained in \texttt{implementation-note.tex} as a historical methods note. This revision organizes the paper around scientific comparison rather than software architecture.

\section{Introduction and research questions}
A constitutive modeller must decide whether an added mechanism improves the predictions required for an application and whether the necessary calibration and initial-state information can be obtained. Counting constants alone does not answer either question. Additional variables can describe memory, density, anisotropy or rate effects, but their usefulness depends on loading paths, state preparation, available observations and numerical implementation.

The intended study combines a broad, source-specific survey with a smaller comparable benchmark. Its core questions are: (i) which mechanisms, state descriptions and calibration requirements distinguish the formulations; (ii) when does additional complexity improve prediction under controlled information; and (iii) which calibration experiments improve transfer? Reviewed parent--extension pairs can address whether a mechanism improves its intended response without degrading other predictions. Behavioural disagreement may subsequently identify useful distinguishing experiments.

The atlas supplies traceable equations and common computational interfaces. This infrastructure is a means of conducting the study. Existing constitutive infrastructure such as MFront already supports executable material behaviours \cite{mfront}; the paper's contribution must be established through its scientific coverage and controlled evidence. A universal leaderboard is not a primary outcome. The proposed outcome is a task-specific account of accuracy, completion, information requirements and computational effort.

Constitutive benchmarking has substantial precedent. D'Onza et al.\ compared models for unsaturated soils with an undisclosed test \cite{donza2011}. Automatic hypoplastic calibration studies examine optimizer and discrepancy choices and repeated-run parameter dispersion \cite{machacek2022}; calibration of both hypoplasticity and SANISAND is also established \cite{machacek2023}. Zeng et al.\ directly examine calibration-dataset composition and weighting for SANISAND-F \cite{zeng2026}. Consequently, calibration-data value or automatic fitting alone cannot establish the present paper's contribution. Its intended distinction is evidence connecting mechanisms across formulations to task-specific prediction under explicit calibration and initialization information. This distinction still requires the completed comparisons and a fuller related-work review.

The closely related fine-sand study by Wichtmann et al.\ compares hypoplasticity with intergranular strain, SANISAND 2004 and ISA 2014 \cite{wichtmann2019}. Its calibration and history choices reinforce the need to account for information beyond fitted curves. Zeng et al.'s full-dataset score includes calibration tests \cite{zeng2026}; it answers a different question from error on disjoint evaluation tests.

\section{Collection, scope and constitutive complexity}
\subsection{Units of inclusion and evidence}
Publications, formulations, implementations and calibrated material parameter sets are distinct records. A parent model, an extension and a different discretization must retain those relationships. The planned survey will document search sources, inclusion decisions, exact variants, source locators and unresolved cases. Final survey counts require completed screening and formulation reconciliation. An accessible implementation or a passing test does not establish independent equation review or experimental validity.

The primary survey covers soils, granular media, intact and jointed rock, and their interfaces. Standalone bulk concrete and geopolymer laws are retained as linked theoretical background. Interfaces with a concrete counterface and cemented or bonded geomaterials remain eligible for source assessment.

A bounded development search uses eight thematic queries, one page of relevance-ordered Crossref results per theme, and a publication-date cutoff of 30 October 2026. It returned \SurveyOccurrenceCount{} occurrences forming \SurveyPublicationCount{} normalized DOI groups. The assisted title pass records \SurveyTitleDecisionCount{} decisions: \SurveyRetrieveCount{} retrieval candidates, \SurveyBackgroundCount{} background/source-tracing routes, \SurveyUnresolvedCount{} unresolved records and \SurveyOutsideScopeCount{} provisional material-scope exclusions; \SurveyUntitledCount{} groups lack a title decision. Exact requests, response hashes and item locators are retained. Independent screening review, primary-source assessment, supplementary discovery and the final search depth remain open. These bibliographic counts are neither formulation counts nor estimates of model prevalence. Search themes do not assign constitutive families.

Survey inclusion and benchmark admission are separate decisions. An annotated formulation may be excluded from a benchmark because of material domain, kinematic restrictions, missing initialization evidence, absent controls or incomplete verification. The eventual benchmark should use a coherent material and loading panel rather than require every surveyed formulation to perform every task.

\subsection{Complexity as several reported quantities}
Report declared input entries, material-role entries, independently adjustable combinations, fitted quantities, supplied prior information and specimen-specific initial/history quantities separately. Independently adjustable material dimension requires a parameterization audit; practical identifiability additionally requires evidence from the actual observations. Inequality constraints, fixed source values and normalization choices must remain explicit.

Stored stress, persistent internal state, bookkeeping and transient integration multipliers are also different quantities. State counts require a kinematic scope and a component convention. An algebraic stress output need not occupy persistent storage; a six-component output under axisymmetric constraints need not have six independently variable components. Unsupported three-dimensional counts remain not applicable rather than zero.

Table~\ref{tab:annotations} describes \AnnotatedFormulationCount{} assisted annotation snapshots, not a screened survey population or admitted benchmark collection. Source hashes and review status accompany each record. Material-role counts in this table are provisional classifications, not identified parameter dimensions.
\begin{table}[htbp]
\centering\small
\caption{Development annotation snapshots. RI denotes rate independent; Axi. denotes the fixed-axis axisymmetric scope. Stored stress and internal columns count components in each supported scope. Bookkeeping and transient multipliers are recorded separately in the source annotations and are excluded from these two columns.}
\label{tab:annotations}
\begin{tabular}{l l r r r r}
\toprule
Model & Scope & \shortstack{Declared\\entries} & \shortstack{Material\\roles} & \shortstack{Stored\\stress} & Internal\\
\midrule
\inputtablerows generated/survey_table.tex
\bottomrule
\end{tabular}
\end{table}

The restricted Figure~1 Terracotta transcription \cite{terracotta} illustrates both distinctions. It stores four evolving scalars and computes stress from state and imposed rate. A coordinated temperature reparameterization identifies a normalization redundancy (Appendix~\ref{app:normalization}). Classifying nine material-role entries after fixing that convention does not prove nine identifiable material combinations. The clay formulation is excluded from the sand pilot and its rate-independent adapter; it is a representation case, not another participant in that comparison.

\subsection{Source examples of calibration burden}
Source review can precede executable transcription. The following examples are partial assisted reviews, separate from the implementation snapshots in Table~\ref{tab:annotations}; neither is admitted to the present benchmark.

The original Nor-Sand example is rigid-plastic and developed for axisymmetric drained loading \cite{norsand1993}. Its Erksak calibration adjusts the hardening coefficient to initial shear stiffness after supplying critical-state and dilatancy information. Density and initial pressure also enter initialization. Counting only the coefficient adjusted in that example would conceal the information already supplied. Later elastic or three-dimensional variants require separate records.

The high-cycle accumulation model advances with cycle count and uses a recording cycle to obtain strain-amplitude information \cite{hca2005}. Its centrifuge example supplies constants for both the accumulation law and a companion hypoplastic model; some companion constants are adjusted using the second loading cycle. The source-case tables therefore describe a combined prediction workflow, rather than establish the independent dimension of one constitutive law. Reference conventions, shared inputs and history initialization require further audit.

These cases motivate reporting calibration burden by prediction task. Accumulation conditional on an observed local cycle and accumulation using cycles predicted by a companion model answer different questions. The former supplies an observation; the latter includes companion-model calibration and amplitude error. Any specimen-specific adaptation must be declared, with subsequent evaluation restricted to the prespecified responses not used for that adaptation. Source-only annotations retain these distinctions while executable representations evolve.

\section{Comparison methods}
\subsection{Prediction tasks and available observations}
Before comparing formulations, define the response scale, controls, conditioning observations, prediction window and scored channels. Keep the formulation's domain separate from the current engine's capabilities. Table~\ref{tab:taskcontract} gives the information declarations needed for three distinct task types; it does not admit those tasks or models to the current benchmark.

\begin{table}[htbp]
\centering\small
\caption{Proposed comparison tasks and information declarations. Geometry, preparation and externally supplied relations belong to the information budget wherever used.}
\label{tab:taskcontract}
\renewcommand{\arraystretch}{1.15}
\begin{tabular}{>{\raggedright\arraybackslash}p{.20\linewidth} >{\raggedright\arraybackslash}p{.34\linewidth} >{\raggedright\arraybackslash}p{.35\linewidth}}
\toprule
Task & Declare before prediction & Evaluate separately\\
\midrule
Homogeneous material point & Imposed stress/strain path, drainage, initial state and history & Prespecified stress, strain or pore-pressure channels under those controls\\
Interface or spatial problem & Geometry, normal constraint, field boundary conditions and any length or thickness information & Traction/slip, displacement and spatial observations supported by the task\\
Combined prediction workflow & Component laws, shared observations, companion calculations and permitted adaptation & Conditional component response versus prediction by the full workflow\\
\bottomrule
\end{tabular}
\end{table}

For each proposed mechanism contrast, establish that at least one available observation responds to the changed ingredient under the actual controls. Equal errors on an insensitive channel cannot establish that an ingredient is unnecessary. Conversely, an internal state variable is not an experimental observation unless a reviewed measurement or inference map connects them. Spatial effects require appropriate boundary-value evidence; homogeneous response alone cannot establish their predictive benefit. These requirements guide task selection while the model schema evolves.

\subsection{Admission and initialization}
Before a controlled study, pin each implementation and establish source-equation review, domain checks, analytical limits, appropriate reference comparisons and numerical accuracy for the proposed tasks. An extension should enter a parent comparison only when both variants pass the relevant admission checks. A common engine reduces some implementation differences but does not eliminate solver effects or establish equal numerical accuracy.

Initialization is part of the information budget. Record measured, supplied, inferred and fitted initial quantities distinctly, including preconsolidation, fabric or history states where required. A changed parameter set may require reinitializing an algebraic-stress model to preserve a prescribed initial stress. Physical-time models additionally require durations, rates and any initial thermal or mesoscopic quantities appropriate to their formulation.

\subsection{Experiments, transformations and whole-test splits}
The development panel draws on the Karlsruhe fine sand database \cite{wt1,wt2}. Raw archive and member hashes, original headers, unit transformations and normalized-record identities are retained. Nominal test-program values do not replace measured initial stress. Specimen preparation, consolidation, geometry, repeatability and grouping evidence must be audited before final statistical claims. Related tests are not automatically independent because their file names differ.

Development tests used for protocol design remain development data. Additional final-evaluation responses were reserved before extracting new development members, with prior exposure to published figures and parameter sets disclosed. No genuinely blind evaluation claim is made. A final protocol must freeze the eligible models, information supplied, bounds, algorithms, metrics and analyses before releasing those reserved responses.

An initial-state conflict remains unresolved: selected file void ratios disagree with the original published table, while a later primary study reports different values \cite{yeh}. The current development choice preserves file initial stress and void ratio; a separate fixed-parameter diagnostic retains the original-table interpretation as a sensitivity scenario. Source selection is not performed by choosing the lower prediction error.

\subsection{Calibration information and selection}
The development pilot includes an elastic diagnostic baseline, the stated perfect-plastic Drucker--Prager specialization, and a partial fit of the von Wolffersdorff hypoplastic law \cite{vw96}. The hypoplastic fit retains source-supplied constants \cite{overshoot}; this extra information prevents an equal-information cross-model complexity interpretation. The intergranular-strain extension \cite{nh97} is annotated separately and is not admitted to this pilot.

The \texttt{pressure3} package contains TMD2, TMD3 and TMD4. The \texttt{density3} package contains TMD3, TMD8 and TMD13. The former spans pressure within the loosest density band; the latter spans density at nominally equal pressure. Both use three tests. They predict the same \PilotEvaluationTestCount{} other development tests: TMD1/5, TMD7/9 and TMD12/14, grouped descriptively by density band. Equal test count does not establish equal preparation cost or equal parameter information. This design was recorded after inspecting the earlier nested development pilot, so it is exploratory.

Calibration uses monotonic ramps through the first \PilotMaximumStrainPercent\% axial strain and \PilotCalibrationSamples{} equally spaced comparison coordinates. For test i and channel c, let \(s_{q,i}=p'_{0,i}\) and \(s_v=\PilotVolumeScale\). The objective is
\[
J_{\rm cal}(\boldsymbol\theta)=\frac{1}{|\mathcal C|}\sum_{i\in\mathcal C}\frac12
\left(\frac{{\rm MSE}_{q,i}}{s_{q,i}^{2}}+
\frac{{\rm MSE}_{v,i}}{s_v^{2}}\right).
\]
The scales are declared physical references, not estimates of measurement noise. Bounds and fixed/fitted partitions are configured before fitting. Bounded coordinate pattern search uses \PilotStartCount{} declared starts and two parameter-vector budgets. The lowest calibration objective selects the fit independently for each package; evaluation error does not select parameters. Search termination and variability do not establish global optimality, posterior uncertainty or identifiability.

\subsection{Acquisition-order evaluation and failure accounting}
All matched-count evaluation tests use the same recorded-order drained replay. The native driver holds radial effective stress while imposing consecutive axial-strain coordinates, including the small TMD12 reversal. Repeated coordinates imply no inferred duration or evolution in these rate-independent laws. This is a prescribed-control interpretation of the recorded strain, not an errors-in-variables treatment or a reconstruction of experimental timing.

Let \(r_{j,c}\) be the endpoint prediction residual and
\(d_j=|\varepsilon_{a,j}-\varepsilon_{a,j-1}|\). The acquisition-order channel error is
\[
E_{c,i}=\frac{\sum_{j=1}^{n_i-1}d_j(r_{j-1,c}^{2}+r_{j,c}^{2})}
{2\sum_{j=1}^{n_i-1}d_j}.
\]
Zero-travel segments receive zero weight. For stress weight w the aggregate is
\[
R_w=\left[\frac{1}{|\mathcal E|}\sum_{i\in\mathcal E}
\left(w\frac{E_{q,i}}{s_{q,i}^{2}}+(1-w)\frac{E_{v,i}}{s_v^{2}}\right)\right]^{1/2}.
\]
Primary evaluation uses w=0.5; w=0.25 and 0.75 are declared evaluation-only sensitivity checks with fitted parameters fixed. This endpoint quadrature differs from the monotonic calibration resampling. Aggregate values must not be compared directly across those metric definitions.

Require complete coverage of each scheduled evaluation window. A failed or missing test prevents a complete aggregate; incomplete cases are retained with their status. Report outside-scope tasks, initialization failures, calibration failures and attempted integration failures separately. A search penalty for failed calibration vectors is not an evaluation ranking rule. Per-test physical errors and completion accompany aggregate errors.

\subsection{Compression initialization and strain coordinates}
Complementary compression data can constrain information inherited by the current partial hypoplastic fit. Before adding those data to calibration, a separate diagnostic uses two development oedometer tests, OE1 and OE12. The remaining ten OE responses were reserved before raw-response inspection. The selected files supply engineering axial strain, axial stress and void ratio, but not lateral stress or a full initial stress tensor. The diagnostic uses the reported axial stress as effective stress under a dry-test interpretation, which remains a qualification for source review. Published simulations start at 1 kPa axial stress and assume vertically mobilized intergranular strain \cite{wt1,wichtmann2019}; those passages do not provide a complete initialization recipe for the present base law.

At fixed material constants, compare assumed initial radial/axial stress ratios \CompressionStressRatios{}, engineering-increment and logarithmic-height mappings, and two maximum numerical increments, \CompressionCoarseIncrement{} and \CompressionFineIncrement{}. Ratios specify initial stress only; radial stress subsequently evolves under zero lateral strain. Each of \CompressionStartThresholdCount{} start thresholds uses the first eligible measured point, without interpolation, and ends at first unloading. This avoids passing zero-stress unloading through a positive-pressure law, but does not reconstruct its full cycle or history. The lowest threshold, \CompressionStartThresholdKPa{} kPa, selects an observed axial stress of \CompressionObservedStartKPa{} kPa in both files. The other starts are development sensitivity choices rather than independently established admissibility limits.

For fixed lateral dimensions and solid volume, engineering strain \(\varepsilon_a\) gives
\[
\frac{1+e_{\rm end}}{1+e_{\rm start}}=
\frac{1-\varepsilon_{a,\rm end}}{1-\varepsilon_{a,\rm start}}.
\]
The native void-ratio continuity law instead integrates to \(1+e_{\rm end}=(1+e_{\rm start})\exp(-\Delta\xi)\) under imposed axial coordinate \(\xi\). Engineering increments and logarithmic height increments therefore supply different coordinates. The latter preserve the stated volume geometry; this mapping alone does not establish finite-strain validity of an infinitesimal stress law. Recorded strain and void ratio are kinematically linked and must not count as independent calibration information. Printed decimal precision is not measurement uncertainty.

\subsection{Numerical and provenance controls}
Selected evaluations use successively smaller maximum increments. State commits are transactional: native trial history is not accumulated during control solves, and a failed replay segment retains the preceding observation state. The replay wrapper checks radial-stress drift against the original reference. Refinement tests numerical sensitivity for the selected histories; it does not independently verify a constitutive transcription.

Each run archives the exact executable, model definitions, effective configuration, experiment copies, paper code, checksums, optimizer journal, start-level fits and termination records. Snapshots protect the scientific identity of a run while atlas sources evolve. The historical pilot uses the previous regularized Drucker--Prager implementation; the current atlas adds an apex branch. These revisions are not merged under an unqualified result label. Concurrent development runtimes do not provide controlled computational-cost comparisons.

\section{Development evidence and its interpretation}
\subsection{Matched calibration-test allocation}
Table~\ref{tab:pilot} and Figure~\ref{fig:weights} are generated from frozen development records. They compare packages within each of \PilotModelCount{} formulations and are not a cross-model leaderboard. Negative paired changes favor density-spanning calibration for the declared evaluation panel.
\begin{table}[!htb]
\centering\small
\caption{Development matched-count comparison at primary evaluation weight. Error is root mean scaled squared error over all scheduled tests. Fits are selected by calibration objective only. Budgets count parameter-vector evaluations per start. The Drucker--Prager rows refer to the historical frozen implementation.}
\label{tab:pilot}
\begin{tabular}{r l r r r}
\toprule
Budget & Formulation & Pressure3 & Density3 & Change\\
\midrule
\inputtablerows generated/pilot_table.tex
\bottomrule
\end{tabular}
\end{table}
\begin{figure}[!htb]
\centering
\IfFileExists{generated/matched-design.png}{\includegraphics[width=\linewidth]{generated/matched-design.png}}{\fbox{Generated development figure pending}}
\caption{Evaluation-weight sensitivity with fitted parameters fixed. Negative changes favor density3. Lines identify search budgets; coincident elastic lines overlap. Panels have different vertical scales. This examines aggregation sensitivity, not recalibration or statistical uncertainty.}
\label{fig:weights}
\end{figure}

The calibration-selected partial hypoplastic fits favor density coverage in the full aggregate at both budgets and the declared weights. Gains concentrate in the higher-density cohort, while the other two cohorts have slightly greater primary error. Same-seed aggregate comparisons retain the direction, but these starts are not independent observations or posterior draws. The selected Drucker--Prager fits also favor density coverage overall, while some individual starts give the opposite preference. The elastic baseline favors pressure coverage in the full aggregate but density coverage in its higher-density cohort. Detailed per-test and start-level records remain in the generated paired report.

These development observations motivate task-specific information analysis. Prior constants, partial calibration, parameter bounds and search performance prevent conclusions about a preferred laboratory programme, mechanism benefit or transfer to other materials.

\subsection{Compression sensitivity before calibration}
All \CompressionReplayCount{} scheduled compression replays completed. Material constants remain the inherited baseline values, rather than the shear pilot's selected fits. The elastic constants are diagnostic choices, and the hypoplastic constants already incorporate source compression information. These runs neither reproduce the complete published IGS case nor provide independent compression validation.

Table~\ref{tab:compression} separates initial-stress, coordinate and numerical effects at the lowest declared start threshold. Comparisons pair identical acquired observations while holding the other factors fixed. The base hypoplastic response is sensitive to both assumed initial stress and coordinate mapping by amounts much larger than the observed refinement differences. The elastic baseline gives the expected zero axial-stress sensitivity to initial lateral stress under its prescribed-strain controls. These checks support explicit treatment of initialization and conventions before interpreting changes from added calibration data. They do not identify a physically correct initial stress, select a coordinate by prediction error, or prove numerical convergence. Different start thresholds use different initial states and intervals and are not a paired common-interval accuracy contrast.

\begin{table}[htbp]
\centering\small
\caption{Fixed-parameter first-loading development sensitivity at the lowest start threshold. Entries are maximum absolute axial-stress differences over matched observation rows. Initial-stress and coordinate columns use the finer numerical resolution. Numerical differences compare the two resolutions at fixed initial stress and mapping. Units differ deliberately to show their magnitudes; these are sensitivities, not prediction errors or uncertainty intervals.}
\label{tab:compression}
\begin{tabular}{l l r r r}
\toprule
Test & Baseline & \shortstack{Initial stress\\(kPa)} & \shortstack{Coordinate\\(kPa)} & \shortstack{Refinement\\(Pa)}\\
\midrule
\inputtablerows generated/compression_table.tex
\bottomrule
\end{tabular}
\end{table}

\subsection{Evidence still required for the main study}
The broad survey, comparable independently reviewed subset, controlled information packages and final protocol remain incomplete. Parent--extension results require admissible reviewed pairs. Physical-time comparisons require rate/duration metadata and appropriate controls. Further tasks include initialization review, calibration-objective sensitivity, stronger optimization checks, parameter/prediction uncertainty, source-matched numerical evidence and broader materials/loading paths. The eventual paper should report computational cost only under declared hardware, builds and verified numerical accuracy.

\section{Discussion and intended conclusions}
The paper should answer what a modeller gains from a mechanism, what measurements make that gain accessible, and where the resulting prediction remains unreliable. The development work already shows why declared parameter count, stored-state size and a single aggregate error cannot substitute for those questions. It also shows why source and implementation identity must remain attached to each result.

Final conclusions will be limited to the reviewed collection and frozen experimental tasks. A material-point comparison does not establish boundary-value predictive accuracy. Related model variants are not independent samples; recording points are not experimental replicates. Until the remaining evidence gates are met, this manuscript presents a study framework and development findings rather than the planned survey and confirmatory benchmark conclusions.

\section*{Data and software availability}
Paper code, study records and manuscript sources are under \texttt{paper/}. Generated manuscript inputs include source-report and output hashes. Local experiment and run snapshots remain excluded from public source distribution pending the applicable data-release review. The archive and normalized source identities are retained for recovery; no DOI or archival-release claim is made. The generator rejects mismatches between paired summaries and their frozen source reports.

\appendix
\section{Normalization and state-count audit}\label{app:normalization}
For the pinned Figure~1 equations, the positive transformation
\[
T_m^*=kT_m,\quad\alpha^*=k\alpha,\quad\gamma^*=k\gamma,
\quad\eta^*=\eta/k,\quad\Gamma^*=k^2\Gamma
\]
leaves the elastic-state evolution and stress observables unchanged when initial temperature is transformed consistently. Temperature evolution is covariant. Equivalently use \(S=T_m/\sqrt\Gamma\), \(A=\alpha/\sqrt\Gamma\), \(C=\gamma/\sqrt\Gamma\) and \(H=\eta\sqrt\Gamma\), giving \(\dot S=A D_v^2+C D_s^2-HS^2\). This paper-side derivation identifies one redundancy and an upper bound of nine material combinations, not a proof of full identifiability. Changing Gamma alone is not this transformation and changes the stress response. Interpreter tests check nonzero-temperature states, signed shear, several rates and positive scale factors. Independent review is pending.

\begin{thebibliography}{99}
\small
\bibitem{mfront} TFEL/MFront developers. Mechanical behaviours in MFront. \url{https://thelfer.github.io/tfel/web/behaviours.html}.
\bibitem{donza2011} F. D'Onza et al. (2011). Benchmark of constitutive models for unsaturated soils. \emph{G\'eotechnique}, 61, 283--302. \url{https://doi.org/10.1680/geot.2011.61.4.283}.
\bibitem{machacek2022} J. Macha\v{c}ek et al. (2022). On the automatic parameter calibration of a hypoplastic soil model. \emph{Acta Geotechnica}, 17, 5253--5273. \url{https://doi.org/10.1007/s11440-022-01669-4}.
\bibitem{machacek2023} J. Macha\v{c}ek, S. Siegel, P. Staubach and H. Zachert (2023). Automatic Parameter Calibration of Two Advanced Constitutive Models. \emph{Lecture Notes in Civil Engineering}, 288, 110--117. \url{https://doi.org/10.1007/978-3-031-12851-6_14}.
\bibitem{zeng2026} S. Zeng, J. Macha\v{c}ek, S. Moshfeghi and M. Taiebat (2026). Data strategies for automatic calibration of soil constitutive models. \emph{Computers and Geotechnics}, 194, 108009. \url{https://doi.org/10.1016/j.compgeo.2026.108009}.
\bibitem{wichtmann2019} T. Wichtmann, W. Fuentes and Th. Triantafyllidis (2019). Inspection of three sophisticated constitutive models based on monotonic and cyclic tests on fine sand: Hypoplasticity vs. Sanisand vs. ISA. \emph{Soil Dynamics and Earthquake Engineering}, 124, 172--183. \url{https://doi.org/10.1016/j.soildyn.2019.05.001}.
\bibitem{vw96} P.-A. von Wolffersdorff (1996). A hypoplastic relation for granular materials with a predefined limit state surface. \emph{Mechanics of Cohesive-Frictional Materials}, 1, 251--271.
\bibitem{nh97} A. Niemunis and I. Herle (1997). Hypoplastic model for cohesionless soils with elastic strain range. \emph{Mechanics of Cohesive-Frictional Materials}, 2, 279--299.
\bibitem{norsand1993} M.~G. Jefferies (1993). Nor-Sand: a simple critical state model for sand. \emph{G\'eotechnique}, 43, 91--103. \url{https://doi.org/10.1680/geot.1993.43.1.91}.
\bibitem{hca2005} A. Niemunis, T. Wichtmann and Th. Triantafyllidis (2005). A high-cycle accumulation model for sand. \emph{Computers and Geotechnics}, 32, 245--263. \url{https://doi.org/10.1016/j.compgeo.2005.03.002}.
\bibitem{wt1} T. Wichtmann and Th. Triantafyllidis (2016). An experimental data base for the development, calibration and verification of constitutive models for sand with focus to cyclic loading. Part I: tests with monotonic loading and stress cycles. \emph{Acta Geotechnica}, 11, 739--761. \url{https://doi.org/10.1007/s11440-015-0402-z}.
\bibitem{wt2} T. Wichtmann and Th. Triantafyllidis (2016). Part II: tests with strain cycles and combined loading. \emph{Acta Geotechnica}, 11, 763--774. \url{https://doi.org/10.1007/s11440-015-0412-x}.
\bibitem{overshoot} J. Macha\v{c}ek, M. Tafili and C. Grandas Tavera (2026). Overshooting of hypoplastic models with intergranular strain: manifestation, implications and remedy. \emph{International Journal for Numerical and Analytical Methods in Geomechanics}, 50, 5068--5091. \url{https://doi.org/10.1002/nag.70375}.
\bibitem{yeh} F.-H. Yeh, M. Tafili, L.~F. Prada-Sarmiento, L. Ge and T. Wichtmann (2023). Inspection of two sophisticated models for sand based on generalized plasticity: Monotonic loading and Monte Carlo analysis. \emph{International Journal for Numerical and Analytical Methods in Geomechanics}, 47, 425--456. \url{https://doi.org/10.1002/nag.3476}.
\bibitem{terracotta} M. Wiebicke and I. Einav (2024). A simple hydrodynamic model for clay. \emph{Journal of the Mechanics and Physics of Solids}, 192, 105789. \url{https://doi.org/10.1016/j.jmps.2024.105789}.
\end{thebibliography}
\end{document}
